\chapter{ŞEKİL VE TABLOLAR}
\section{Şekil Atıflar ve Şekil Örneği}
Tablo ve Şekiller sayfa düzeni esaslarına uymak koşuluyla, metinde ilk söz edildikleri yerden hemen sonraya veya mümkün olduğu ölçüde en yakınına yerleştirilmelidir. Tablo ve şekillerden önce, ilgili tablo ya da şekile atıfta bulunulmalıdır.

Şekil başlığı şeklin altına, tablo başlığı ise tablonun üstüne yazılır. Tablo ve şekil başlıkları bir satıra sığıyorsa yazı bloğuna göre ortalı olarak yazılır. Bir satırdan uzun olan başlıklar ise iki yana yaslı olarak yazılmalıdır. Bu durumda, ikinci satır, tablo ve şekil numarasından sonra ilk satır metni ile aynı hizadan başlayacak şekilde girintili olmalıdır.

Tablo ve şekillere, ilk rakam bölüm numarası (Ekler bölümünde ise harf), ikinci rakam tablo veya şeklin bölüm içindeki sıra numarası olmak üzere numara verilir. Tablo, şekil kelimeleri ve numaraları koyu harflerle yazılır ve sonlarına nokta konulmalıdır.  (Örnek: Tablo 1.2., Şekil 3.5., Tablo A.1., Şekil B.5.).

Tablo ve şekillere metin içerisinde atıf yapılırken tablo veya şekil numaralarının sonundaki nokta kullanılmaz. (Örneğin; Şekil 1.1'de görüldüğü gibi \ldots{})

Tablo başlıkları, bir satır aralıklı olarak, başlıktan önce 12 punto (nk), sonra ise 6 punto (nk) aralık bırakılarak yazılmalı ve Tablo başlıkları nokta ile bitirilmelidir. Tablo başlığı ile tablo farklı sayfalarda yer alamaz.

Tablodan sonra gelen metin bölümündeki ilk paragraf üstten 12 punto (nk) ve alttan 6 punto (nk) aralık bırakılarak yazılmalıdır. Tablolardan hemen sonra gelecek başlıklar, belirtilen başlık formatlarında değişiklik yapılmadan aynen kullanılmalıdır.

Tablolarda kullanılacak dipnotlar (gerekliyse) 1 satır aralıklı ve metinden 2 yazı boyutu küçük olarak yazılmalıdır.

Şekil başlıkları, bir satır aralıklı olarak, başlıktan önce 6 punto (nk), sonra ise 12 punto (nk) aralık bırakılarak yazılmalı ve şekil başlıkları nokta ile bitirilmelidir. Şekil başlığı ile şekil farklı sayfalarda yer alamaz. Şekilden önce gelen metin bölümündeki son paragraf, önce 6 punto (nk), sonra 12 punto (nk) aralık bırakılarak yazılmalıdır.

Bir sayfayı aşan büyüklükteki tablo ve şekillerde 2. sayfada aynı tablo/şekil numarası ve açıklaması yazılarak, tablo/şekil numarası ile açıklaması arasına, parantez içinde (Devamı) yazılmalıdır. (Örneğin; Tablo 1.1. (Devamı):\textbf{ }Geometrik yapı parametreleri, Şekil 1.1. (Devamı):\textbf{ }Elektronik soğurma spektrumu).

Tezlerde kullanılacak olan grafik ve resimler şekil olarak dikkate alınmalı ve buna göre düzenlenmelidir.

Birden fazla Tablo veya Şekil aynı sayfada yer alabilir. Ancak, 4 sayfadan daha uzun olan tablo veya şekiller ek olarak verilmelidir.

% [Normal] sp=after:240
Katlı sayfa (A4 ebadından büyük olan) ve sayfa üzerine iliştirilmiş görsel malzeme gibi sayfa kalınlığını arttırarak tezin açılma düzenini bozan sayfalar da ek olarak verilmelidir.

\begin{figure}[H]
  \centering
  \includegraphics[width=11.11cm]{sekiller/image1.png}
  \caption{Tüm şekil ve tablolar ile bunların başlıkları yazı bloğuna göre ortalı olarak yerleştirilmelidir.}
\end{figure}
Tablo ve şekillerde gerekli ise yazı boyutu 8 punto yazı boyutuna kadar küçültülebilir.

Tablolar tezde kullanılan yazı karakteriyle hazırlanır, şekillerde kullanılan yazı karakteri tez boyunca kendi içerisinde tutarlı olmalıdır.

% [Normal] sp=after:240
Tablolar ve şekiller sayfa düzeni esaslarına uymak şartı ile metinde ilk söz edildikleri yerden hemen sonraya mümkün olduğu kadar yakın yerleştirilmelidir (Şekil 2.1.). Tablo ve şekillerden önce, ilgili tablo ya da şekile atıfta bulunulmalıdır (Tablo 1.1).

\begin{figure}[H]
  \centering
  \includegraphics[width=11.11cm]{sekiller/image1.png}
  \caption{Bir satıra sığan şekil başlıkları ortalı olarak yerleştirilmelidir.}
\end{figure}
Şekil 2.2’de ki Sed diam nonumy eirmod tempor invidunt ut labore et dolore magna aliquyam erat, sed diam voluptua. At vero eos et accusam et justo duo dolores et ea rebum. Lorem ipsum dolor sit amet, consetetur sadipscing elitr, sed diam nonumy eirmod tempor invidunt ut labore et dolore magna aliquyam erat, sed diam voluptua. At vero eos et accusam et justo duo dolores et ea rebum. At vero eos et accusam et justo duo dolores et ea rebum. At vero eos et accusam et justo duo dolores et ea rebum.

Lorem ipsum dolor sit amet, consetetur sadipscing elitr, sed diam nonumy eirmod tempor invidunt ut labore et dolore magna aliquyam erat, sed diam voluptua. At vero eos et accusam et justo duo dolores et ea rebum. At vero eos et accusam et justo duo dolores et ea rebum. At vero eos et accusam et justo duo dolores et ea rebum.

\section{Yatay Sayfada Şekil Örneği}
Lorem ipsum dolor sit amet, consetetur sadipscing elitr, sed diam nonumy eirmod tempor invidunt ut labore et dolore magna aliquyam erat, sed diam voluptua. At vero eos et accusam et justo duo dolores et ea rebum (Şekil 2.3). Lorem ipsum dolor sit amet, consetetur sadipscing elitr, sed diam nonumy eirmod tempor invidunt ut labore et dolore magna aliquyam erat, sed diam voluptua. At vero eos et accusam et justo duo dolores et ea rebum.

Lorem ipsum dolor sit amet, consetetur sadipscing elitr, sed diam nonumy eirmod tempor invidunt ut labore et dolore magna aliquyam erat, sed diam voluptua. At vero eos et accusam et justo duo dolores et ea rebum. Lorem ipsum dolor sit amet, consetetur sadipscing elitr, sed diam nonumy eirmod tempor invidunt ut labore et dolore magna aliquyam erat, sed diam voluptua. At vero eos et accusam et justo duo dolores et ea rebum. Lorem ipsum dolor sit amet, consetetur sadipscing elitr, sed diam nonumy eirmod tempor invidunt ut labore et dolore magna aliquyam erat, sed diam voluptua. At vero eos et accusam et justo duo dolores et ea rebum.

\begin{landscape}
\begin{figure}[H]
  \centering
  \includegraphics[width=21.96cm,height=10.54cm]{sekiller/image2.png}
  \caption{Yatay tam sayfa şekil.}
\end{figure}
\end{landscape}
Lorem ipsum dolor sit amet, consetetur sadipscing elitr, sed diam nonumy eirmod tempor invidunt ut labore et dolore magna aliquyam erat, sed diam voluptua. At vero eos et accusam et justo duo dolores et ea rebum. Lorem ipsum dolor sit amet, consetetur sadipscing elitr, sed diam nonumy eirmod tempor invidunt ut labore et dolore magna aliquyam erat, sed diam voluptua. At vero eos et accusam et justo duo dolores et ea rebum. Lorem ipsum dolor sit amet, consetetur sadipscing elitr, sed diam nonumy eirmod tempor invidunt ut labore et dolore magna aliquyam erat, sed diam voluptua. At vero eos et accusam et justo duo dolores et ea rebum. Lorem ipsum dolor sit amet, consetetur sadipscing elitr, sed diam nonumy eirmod tempor invidunt ut labore et dolore magna aliquyam erat, sed diam voluptua. At vero eos et accusam et justo duo dolores et ea rebum.

\section{Tablo Atıfları ve Tablo Örneği}
Lorem ipsum dolor sit amet, consetetur sadipscing elitr, sed diam nonumy eirmod tempor invidunt ut labore et dolore magna aliquyam erat, sed diam voluptua. At vero eos et accusam et justo duo dolores et ea rebum. Stet clita kasd gub rgren, no sea takimata sanctus est Lorem ipsum dolor sit amet, consetetur sadipscing elitr, sed diam nonumy eirmod tempor invidunt ut lab ore sit et dolore magna.

Tablolar sayfa düzeni esaslarına uymak koşuluyla metinde ilk söz edildikleri yerden hemen sonraya veya mümkün olduğu ölçüde en yakınına yerleştirilmelidir. Tablolardan önce, ilgili tabloya atıfta bulunulmalıdır.

Her tablonun başlığı tablonun üstüne yazılır. Tüm Tablo başlıkları bir satıra sığıyorsa yazı bloğuna göre ortalı olarak yerleştirilmelidir. Bir satırdan uzun olan başlıklar ise iki yana yaslı olarak yerleştirilmelidir. İkinci satır, tablo ve şekil numarasından sonra ilk satır metni ile aynı hizadan başlayacak şekilde girintili olmalıdır.

Tablolarda ilk rakam bölüm numarası (Ekler bölümünde harf), ikinci rakam tabloların  bölüm içindeki sıra numarası olmak üzere numara verilir. Örneklerde olduğu gibi tablo kelimesi ve numaraları koyu harflerle yazılır ve sonunda nokta konulmalıdır.  (Örnek: Tablo 1.2., Tablo A.1.).

Tablo başlıkları, bir satır aralığı kullanılarak ve yazıdan önce 12 punto (nk), sonra 6 punto (nk)  aralık bırakılarak yazılmalı ve Tablo başlığı nokta ile bitirilmelidir. Tablo başlığı ile tablonun tamamı aynı sayfa içinde yer almalıdır.

Tablodan sonra gelen metin bölümündeki ilk paragraf üstten 12 punto (nk) ve alttan 6 punto (nk) aralık bırakılarak yazılmalıdır. Tablolardan hemen sonra gelecek başlıklar, belirtilen başlık formatlarında değişiklik yapılmadan aynen kullanılmalıdır.

Tablolarda dipnot kullanılması gerekiyorsa 1 satır aralıklı ve metinden 2 yazı boyutu küçük yazılmalıdır.

Bir sayfayı aşan büyüklükteki tablolarda 2. sayfada aynı tablo numarası ve açıklaması yazılarak, tablo numarası ile açıklaması arasına, parantez içinde (Devamı) yazılmalıdır. (Örneğin; Tablo 1.1. (Devamı)\textbf{ }Atıklardaki metal içerikleri).

Birden fazla Tablo aynı sayfada yer alabilir. Ancak 4 sayfadan daha uzun olan tablolar ek olarak verilmelidir.

Katlı sayfa (A4 ebadından büyük olan) ve sayfa üzerine iliştirilmiş görsel malzeme gibi sayfa kalınlığını arttırarak tezin açılma düzenini bozan sayfalar ek olarak verilmelidir.

Tablo 2.1’de görüldüğü üzere Lorem ipsum dolor sit amet, consetetur sadipscing elitr, sed diam nonumy eirmod tempor invidunt ut labore et dolore magna aliquyam erat, sed diam voluptua. At vero eos et accusam et justo duo dolores et ea rebum. Stet clita kasd gub rgren, no sea takimata sanctus est Lorem ipsum dolor sit amet, consetetur sadipscing elitr, sed diam nonumy eirmod tempor invidunt ut lab ore sit et dolore magna.

\begin{table}[H]
  \caption{Tek satırlı ve kolonlar ortalanmış tablo.}
  \centering
\begin{tabular}{C{3.09cm}C{2.13cm}C{2.19cm}C{2.56cm}}
\hline\hline
\Ns Kolon A & Kolon B & Kolon C & Kolon D \\
\hline
\Ns Satır A & Satır A & Satır A & Satır A \\
\Ns Satır B & Satır B & Satır B & Satır B \\
\Ns Satır C & Satır C & Satır C & Satır C \\
\hline
\end{tabular}
\end{table}
% [Normal] sp=before:240
Lorem ipsum dolor sit amet, consetetur sadipscing elitr, sed diam nonumy eirmod tempor invidunt ut labore et dolore magna aliquyam erat, sed diam voluptua. At vero eos et accusam et justo duo dolores et ea rebum.

Lorem ipsum dolor sit amet, consetetur sadipscing elitr, sed diam nonumy eirmod tempor invidunt ut labore et dolore magna aliquyam erat, sed diam voluptua. At vero eos et accusam et justo duo dolores et ea rebum. Stet clita kasd gub rgren, no sea takimata sanctus est Lorem ipsum dolor sit amet, consetetur sadipscing elitr, sed diam nonumy eirmod tempor invidunt ut lab ore sit et dolore magna. Lorem ipsum dolor sit amet, consetetur sadipscing elitr, sed diam nonumy eirmod tempor invidunt ut labore et dolore magna aliquyam erat, sed diam voluptua. At vero eos et accusam et justo duo dolores et ea rebum.

\begin{table}[H]
  \caption{Tablo başlığı nokta ile bitirilmelidir.}
  \centering
\begin{tabular}{C{3.05cm}C{2.10cm}C{2.16cm}C{2.52cm}}
\hline\hline
\Ns Kolon A & Kolon B & Kolon C & Kolon D \\
\hline
\Ns Satır A & Satır A & Satır A & Satır A \\
\Ns Satır B & Satır B & Satır B & Satır B \\
\Ns Satır C & Satır C & Satır C & Satır C \\
\hline
\end{tabular}
\end{table}
% [GOVDE] sp=before:240
Lorem ipsum dolor sit amet, consetetur sadipscing elitr, sed diam nonumy eirmod tempor invidunt ut labore et dolore magna aliquyam erat, sed diam voluptua. At vero eos et accusam et justo duo dolores et ea rebum. Lorem ipsum dolor sit amet, consetetur sadipscing elitr, sed diam nonumy eirmod tempor invidunt ut labore et dolore magna aliquyam erat, sed diam voluptua. At vero eos et accusam et justo duo dolores et ea rebum. Lorem ipsum dolor sit amet, consetetur sadipscing elitr, sed diam nonumy eirmod tempor invidunt ut labore et dolore magna aliquyam erat, sed diam voluptua. At vero eos et accusam et justo duo dolores et ea rebum.

\begin{table}[H]
  \caption{Düşey sayfada devam eden tablo örneği.}
  \centering
\begin{tabular}{L{1.72cm}C{1.80cm}C{2.22cm}R{1.04cm}R{1.04cm}R{1.05cm}}
\hline\hline
\multirow{2}{*}{Parametre} & \multirow{2}{*}{Kolon 2} & \multirow{2}{*}{Kolon 3} & \multicolumn{3}{c}{\Ts Kolon 4} \\
\cline{4-6}
 &  &  & \Ts Alt kolon & Alt kolon & Alt kolon \\
\hline
\Ts Satır 1 & -7.680442 & 7.6986348 & 0.00 & 0.00 & 0.00 \\
\Ts Satır 2 & 140 & - & 0.50 & 0.00 & 0.00 \\
\Ts Satır 3 & 37.174357 & 37.16192697 & 0.00 & 0.00 & 0.00 \\
\Ts Satır 4 & 140 & - & 0.50 & 0.00 & 0.00 \\
\Ts Satır 5 & 37.174357 & 37.16192697 & 0.00 & 0.00 & 0.00 \\
\hline
\end{tabular}
\end{table}
\begin{table}[H]
  \ContinuedFloat
  \caption[]{\devami{} Düşey sayfada devam eden tablo örneği.}
  \centering
\begin{tabular}{L{1.72cm}C{1.80cm}C{2.22cm}R{1.04cm}R{1.04cm}R{1.05cm}}
\hline\hline
\multirow{2}{*}{Parametre} & \multirow{2}{*}{Kolon 2} & \multirow{2}{*}{Kolon 3} & \multicolumn{3}{c}{\Ts Kolon 4} \\
\cline{4-6}
 &  &  & \Ts Alt kolon & Alt kolon & Alt kolon \\
\hline
\Ts Satır 1 & -7.680442 & 7.6986348 & 0.00 & 0.00 & 0.00 \\
\Ts Satır 2 & 140 & - & 0.50 & 0.00 & 0.00 \\
\Ts Satır 3 & 37.174357 & 37.16192697 & 0.00 & 0.00 & 0.00 \\
\Ts Satır 4 & 140 & - & 0.50 & 0.00 & 0.00 \\
\Ts Satır 5 & 37.174357 & 37.16192697 & 0.00 & 0.00 & 0.00 \\
\Ts Satır 6 & 140 & - & 0.50 & 0.00 & 0.00 \\
\Ts Satır 7 & 37.174357 & 37.16192697 & 0.00 & 0.00 & 0.00 \\
\hline
\end{tabular}
\end{table}
% [GOVDE] sp=before:240
Loremipsum dolor sit amet, consetetur sadipscing elitr, sed diam nonumy eirmod tempor invidunt ut labore et dolore magna aliquyam erat, sed diam voluptua. At vero eos et accusam et justo duo dolores et ea rebum. Lorem ipsum dolor sit amet, consetetur sadipscing elitr, sed diam nonumy eirmod tempor invidunt ut labore et dolore magna aliquyam erat, sed diam voluptua. At vero eos et accusam et justo duo dolores et ea rebum. Lorem ipsum dolor sit amet, consetetur sadipscing elitr, sed diam nonumy eirmod tempor invidunt ut labore et dolore magna aliquyam erat, sed diam voluptua. At vero eos et accusam et justo duo dolores et ea rebum.

\section{Yatay Sayfada Tablo Örneği}
% [Normal] sp=before:240
Lorem ipsum dolor sit amet, consetetur sadipscing elitr, sed diam nonumy eirmod tempor invidunt ut labore et dolore magna aliquyam erat, sed diam voluptua. At vero eos et accusam et justo duo dolores et ea rebum. Stet clita kasd gub rgren, no sea takimata sanctus est Lorem ipsum dolor sit amet, consetetur sadipscing elitr, sed diam nonumy eirmod tempor invidunt ut lab ore sit et dolore magna.

\begin{landscape}
\begin{table}[H]
  \caption{2. Satıra geçen örnek tablo adı, 2. Satıra geçen örnek tablo adı, 2. Satıra geçen örnek tablo adı, 2. Satıra geçen örnek tablo adı, 2. Satıra geçen örnek tablo adı.}
  \centering
\begin{tabular}{L{72.6bp}L{76.05bp}L{92.95bp}L{68.65bp}L{68.65bp}L{68.65bp}L{68.65bp}L{68.4bp}}
\hline\hline
\multirow{2}{*}{Parametre} & \multirow{2}{*}{Kolon 2} & \multirow{2}{*}{Kolon 3} & \multicolumn{3}{l|}{\Ts Kolon 4} & \multicolumn{2}{l}{Kolon 5} \\
\cline{4-8}
 &  &  & \Ts Alt kolon & Alt kolon & \multicolumn{1}{l|}{Alt kolon} & Alt kolon & Alt kolon \\
\hline
\Ts Satır 1 & -7.680442 & 7.6986348 & 0.00 & 0.00 & 0.00 & 12 & 12 \\
\Ts Satır 2 & 140 & - & 0.50 & 0.00 & 0.00 & 0 & 0 \\
\Ts Satır 3 & 37.174357 & 37.16192697 & 0.00 & 0.00 & 0.00 & 0 & 24 \\
\Ts Satır 4 & 140 & - & 0.50 & 0.00 & 0.00 & 0 & 0 \\
\Ts Satır 5 & 37.174357 & 37.16192697 & 0.00 & 0.00 & 0.00 & 0 & 24 \\
\Ts Satır 6 & 140 & - & 0.50 & 0.00 & 0.00 & 0 & 0 \\
\Ts Satır 7 & 37.174357 & 37.16192697 & 0.00 & 0.00 & 0.00 & 0 & 24 \\
\Ts Satır 8 & 140 & - & 0.50 & 0.00 & 0.00 & 0 & 0 \\
\Ts Satır 9 & 37.174357 & 37.16192697 & 0.00 & 0.00 & 0.00 & 0 & 24 \\
\Ts Satır 10 & 140 & - & 0.50 & 0.00 & 0.00 & 0 & 0 \\
\Ts Satır 11 & 37.174357 & 37.16192697 & 0.00 & 0.00 & 0.00 & 0 & 24 \\
\Ts Satır 12 & 140 & - & 0.50 & 0.00 & 0.00 & 0 & 0 \\
\Ts Satır 13 & 37.174357 & 37.16192697 & 0.00 & 0.00 & 0.00 & 0 & 24 \\
\Ts Satır 14 & 140 & - & 0.50 & 0.00 & 0.00 & 0 & 0 \\
\Ts Satır 15 & 37.174357 & 37.16192697 & 0.00 & 0.00 & 0.00 & 0 & 24 \\
\hline
\end{tabular}
\end{table}
\begin{table}[H]
  \ContinuedFloat
  \caption[]{\devami{} 2. Satıra geçen örnek tablo adı, 2. Satıra geçen örnek tablo adı, 2. Satıra geçen örnek tablo adı, 2. Satıra geçen örnek tablo adı, 2. Satıra geçen örnek tablo adı.}
  \centering
\begin{tabular}{L{69.0bp}L{72.2bp}L{88.35bp}L{65.15bp}L{65.15bp}L{65.15bp}L{65.15bp}L{64.9bp}}
\hline\hline
\multirow{2}{*}{Parametre} & \multirow{2}{*}{Kolon 2} & \multirow{2}{*}{Kolon 3} & \multicolumn{3}{l|}{\Ts Kolon 4} & \multicolumn{2}{l}{Kolon 5} \\
\cline{4-8}
 &  &  & \Ts Alt kolon & Alt kolon & \multicolumn{1}{l|}{Alt kolon} & Alt kolon & Alt kolon \\
\hline
\Tsd Satır 16 & -7.680442 & 7.6986348 & 0.00 & 0.00 & 0.00 & 12 & 12 \\
\Tsd Satır 17 & 140 & - & 0.50 & 0.00 & 0.00 & 0 & 0 \\
\Tsd Satır 18 & 37.174357 & 37.16192697 & 0.00 & 0.00 & 0.00 & 0 & 24 \\
\Tsd Satır 19 & 140 & - & 0.50 & 0.00 & 0.00 & 0 & 0 \\
\Tsd Satır 20 & 37.174357 & 37.16192697 & 0.00 & 0.00 & 0.00 & 0 & 24 \\
\Tsd Satır 21 & 140 & - & 0.50 & 0.00 & 0.00 & 0 & 0 \\
\Tsd Satır 22 & 37.174357 & 37.16192697 & 0.00 & 0.00 & 0.00 & 0 & 24 \\
\Tsd Satır 23 & 140 & - & 0.50 & 0.00 & 0.00 & 0 & 0 \\
\hline
\end{tabular}
\end{table}

\end{landscape}

